% caveats_d2.tex -- d = 2 caveats (from phi4HT2d, 2026-10-14; readability review 2026-10-15).
\paragraph{Two dimensions.}\label{ht:d2:caveats}
All $d=2$ values are finite-$N$ values of the truncated model. With cutoffs up to $N=24$, and a behaviour that changes
with $\Ds$, we quote no continuum limit and no extrapolation error. The window model is known to be correct only
through third order in $\eps$: expanded in $\eps$, the model agrees with the consistent-window scheme, whose
fourth-order coefficient grows like $-0.027\ln N$ (Section~\ref{ht:d2:weak}). Even at weak coupling, the $N\to\infty$
limit of the truncated model is therefore not the continuum beyond order $\eps^3$. Toward the short-range end the
method fails: at $\Ds\le0.175$ the values fall with growing steps, at $\Ds=0.125$ the value at $N=22$ lies $6\%$ below
the exact $\DE=1$ and moves away, and at $N=24$ the tuning conditions have no solution for $\Ds\lesssim0.1274$. The rise
at $0.25\le\Ds\le0.4$ bounds the truncated model's limit from below only if it continues.

The fold location is an estimate from two values of $\Ds$, and it assumes a saddle-node. The absence of a root at
$N=24$ is established on the scanned range $12.5\le\theta_0\le15.5$, where $f_1$ has a single interior minimum and rises
steeply at both ends. A root on another branch, far from the one continued from smaller $N$, is not excluded. Finally,
the independent solver is not blind (Section~\ref{ht:d2:checks}).
